\section{Results}
\label{sec:results}

% Drafting note:
% This Results outline follows the verified-only artifact path. Keep WordPress
% claims descriptive because the WordPress corpus has no gold Lemotif labels.
% Unsupported planned analyses such as SHAP, longitudinal modeling, weak-label
% training, and inter-annotator reliability should be moved to Methods revisions,
% limitations, or future work unless new artifacts are added.

\subsection{Results Overview}
\label{subs:results-overview}

\begin{itemize}
    \item This chapter reports the completed evaluation pipeline in four steps: dataset distribution analysis, in-domain Lemotif model evaluation, external Stories evaluation, and descriptive WordPress inference.
    \item RQ1 is addressed through the Lemotif held-out test and the human-labeled Stories external test.
    \item RQ2 is addressed through public WordPress journal inference, but only descriptively. Because the WordPress corpus is not human-labeled in the Lemotif taxonomy, this section cannot report WordPress accuracy, $F_1$, or formal bias metrics.
    \item The central narrative is that BERT learns broad affective structure and derived sentiment more reliably than exact 18-label emotion sets, while Stories and WordPress expose domain-shift limitations.
\end{itemize}

\subsection{Dataset and Label Distribution}
\label{subs:dataset-label-distribution}

\begin{itemize}
    \item Begin by explaining that multi-label emotion classification must be interpreted in relation to label prevalence because rare labels strongly affect macro-$F_1$ and per-label results.
    \item Report the Lemotif sample size: 1,473 cleaned reflections.
    \item Report the Lemotif derived sentiment distribution: positive 81.7\% ($n = 1{,}204$), negative 14.7\% ($n = 217$), and neutral 3.5\% ($n = 52$).
    \item Interpret the positive skew using the Pollyanna hypothesis/principle: evaluatively positive language tends to occur more frequently and diversely than evaluatively negative language in ordinary communication \parencite{boucher1969pollyanna,matlin1978pollyanna}.
    \item Add thesis-specific context: the Lemotif prompt asked participants to describe salient aspects of the previous day, which plausibly elicits ordinary-life positive reflection rather than clinical distress narratives.
    \item Note an additional possible skew: crowdsourced short reflections may favor socially acceptable and easily narratable events, while rare negative labels such as Jealous, Afraid, Disgusted, Ashamed, and Awkward occur sparsely.
    \item Compare Stories as a harder external test: positive 65.8\% ($n = 75$), negative 22.8\% ($n = 26$), and neutral 11.4\% ($n = 13$).
    \item End this subsection by previewing that high derived-sentiment performance on Lemotif should not be assumed to transfer directly to Stories.
\end{itemize}

\begin{figure}[htbp]
    \centering
    \includegraphics[width=\textwidth]{figures/results/main_text/fig01_lemotif_emotion_prevalence.png}
    \caption{Emotion label prevalence in the Lemotif dataset. The distribution shows strong positive-affect dominance, which is relevant for interpreting macro-$F_1$ and rare-label performance.}
    \label{fig:lemotif-emotion-prevalence}
\end{figure}

\begin{figure}[htbp]
    \centering
    \includegraphics[width=\textwidth]{figures/results/main_text/fig02_lemotif_stories_distribution_shift.png}
    \caption{Emotion label prevalence across Lemotif and Stories. The Stories test set is less positive and more affectively mixed, indicating a distribution shift between training and external evaluation data.}
    \label{fig:lemotif-stories-distribution-shift}
\end{figure}

% Appendix links to cite in prose:
% Appendix A, Figures A1-A5: word count, topic prevalence, topic-emotion heatmap,
% emotion co-occurrence, and grouped emotion prevalence.

\subsection{In-Domain Emotion Classification on Lemotif}
\label{subs:in-domain-emotion-classification}

\begin{itemize}
    \item Introduce the fine-tuned BERT model as the primary in-domain baseline. BERT is a transformer language model that can be fine-tuned for downstream classification with limited task-specific architecture changes \parencite{vaswani2017attention,devlin2019bert}.
    \item Report the held-out Lemotif test size: 221 reflections.
    \item Report headline emotion metrics: micro-$F_1 = .502$, macro-$F_1 = .319$, exact subset accuracy $= .009$, and hamming loss $= .177$.
    \item Report derived sentiment metrics: accuracy $= .914$ and macro-$F_1 = .652$.
    \item Interpret micro-$F_1$ and macro-$F_1$ separately: micro-$F_1$ captures average label-level decisions across all labels, while macro-$F_1$ exposes poor performance on sparse labels.
    \item Explain that exact subset accuracy is low because a multi-label prediction is counted as exactly correct only if the full predicted emotion set matches the gold set.
    \item Discuss strongest labels: Happy ($F_1 = .766$, $n = 120$), Frustrated ($F_1 = .750$, $n = 28$), and Satisfied ($F_1 = .607$, $n = 86$).
    \item Discuss weakest nonzero-support labels: Afraid ($F_1 = .000$, $n = 1$), Jealous ($F_1 = .000$, $n = 1$), and Disgusted ($F_1 = .091$, $n = 2$).
    \item Conclude that the model supports in-domain affect summarization, especially sentiment-level summarization, but not deployment-ready exact emotion-set prediction.
\end{itemize}

\begin{table}[htbp]
    \centering
    \caption{In-domain Lemotif test performance for the main BERT emotion classifier.}
    \label{tab:lemotif-bert-main}
    \begin{tabular}{lrrrrrr}
        \toprule
        Model & Micro-$F_1$ & Macro-$F_1$ & Subset acc. & Hamming loss & Sent. acc. & Sent. macro-$F_1$ \\
        \midrule
        BERT baseline & .502 & .319 & .009 & .177 & .914 & .652 \\
        \bottomrule
    \end{tabular}
\end{table}

\begin{figure}[htbp]
    \centering
    \includegraphics[width=\textwidth]{figures/results/main_text/fig03_bert_validation_test_metrics.png}
    \caption{Validation and held-out test metrics for the in-domain BERT emotion classifier.}
    \label{fig:bert-validation-test-metrics}
\end{figure}

\begin{figure}[htbp]
    \centering
    \includegraphics[width=\textwidth]{figures/results/main_text/fig04_bert_per_emotion_f1.png}
    \caption{Per-emotion $F_1$ scores on the Lemotif held-out test set. Frequent labels such as Happy and Satisfied are detected more reliably than sparse labels such as Afraid and Jealous.}
    \label{fig:bert-per-emotion-f1}
\end{figure}

\begin{figure}[htbp]
    \centering
    \includegraphics[width=.82\textwidth]{figures/results/main_text/fig05_sentiment_confusion_matrix.png}
    \caption{Confusion matrix for derived sentiment on the Lemotif held-out test set. Derived sentiment is substantially more reliable than exact multi-label emotion prediction.}
    \label{fig:sentiment-confusion-matrix}
\end{figure}

\subsection{Model Variants and Inference Tradeoffs}
\label{subs:model-variants-inference-tradeoffs}

\begin{itemize}
    \item Present Optuna as a hyperparameter-tuned model variant, not the sole main model. Optuna increases emotion micro-$F_1$ to .566 and improves hamming loss to .127, but macro-$F_1$ remains low at .299.
    \item Present fixed inference as a threshold/cardinality adjustment, not a retrained model. It uses the unchanged fine-tuned BERT weights, threshold .59, and a minimum one-label fallback.
    \item Explain the practical tradeoff: fixed inference reduces overprediction and increases derived sentiment accuracy from .914 to .941, but emotion micro-$F_1$ decreases from .502 to .445.
    \item Report the key density change: average predicted labels decrease from 4.330 to 2.127; rows with zero predicted emotions decrease from 9 to 0; rows with at least two extra predicted emotions decrease from 153 to 27.
    \item Interpret fixed inference as useful for readable affect summaries, not as the strongest fine-grained emotion classifier.
    \item Mention that Optuna was implemented with a hyperparameter optimization framework designed for flexible search and pruning \parencite{akiba2019optuna}.
\end{itemize}

\begin{table}[htbp]
    \centering
    \caption{Comparison of BERT model variants on the Lemotif held-out test set.}
    \label{tab:bert-variants}
    \begin{tabular}{lrrrrrr}
        \toprule
        Model & Micro-$F_1$ & Macro-$F_1$ & Subset acc. & Hamming loss & Sent. acc. & Avg. labels \\
        \midrule
        BERT baseline & .502 & .319 & .009 & .177 & .914 & 4.330 \\
        BERT Optuna & .566 & .299 & .072 & .127 & .923 & 3.172 \\
        BERT fixed inference & .445 & .258 & .109 & .130 & .941 & 2.127 \\
        \bottomrule
    \end{tabular}
\end{table}

\begin{figure}[htbp]
    \centering
    \includegraphics[width=\textwidth]{figures/results/main_text/fig06_fixed_inference_cardinality.png}
    \caption{Prediction-cardinality comparison for true labels, baseline BERT predictions, and fixed-inference BERT predictions. The adjusted inference setting reduces label overestimation but lowers fine-grained emotion $F_1$.}
    \label{fig:fixed-inference-cardinality}
\end{figure}

% Appendix links to cite in prose:
% Appendix B: fixed-inference overall metrics and bootstrap confidence intervals.
% Appendix C: threshold sensitivity and calibration diagnostics.

\subsection{Error Analysis and Interpretability}
\label{subs:error-analysis-interpretability}

\begin{itemize}
    \item Use this section to satisfy the rubric requirement for error-pattern analysis.
    \item Start with the core error pattern: the model performs better on frequent and affectively clear labels than rare or ambiguous labels.
    \item Discuss rare labels as a data limitation as well as a model limitation. Afraid and Jealous each have one held-out example, so stable label-specific evaluation is not possible for those emotions.
    \item Discuss multi-label overlap. Reflections can contain blended emotions, so errors often involve missing one emotion while predicting a semantically adjacent or co-occurring label.
    \item Discuss derived sentiment errors. Broad valence is easier than exact emotion identity, but neutral remains weaker because neutral reflections may contain low-intensity affect or mixed signals.
    \item Describe LIME as a local interpretability method that explains individual predictions by fitting interpretable surrogate models around them \parencite{ribeiro2016why}. Do not present LIME as global causal explanation.
    \item Describe feature-ablation outputs as lightweight diagnostics rather than definitive feature importance. The available feature diagnostics use scikit-learn-style modeling tools \parencite{pedregosa2011scikit}.
    \item Use Appendix C for rare-label errors, multilabel error-flow heatmaps, threshold calibration, feature-ablation diagnostics, and LIME panels.
\end{itemize}

Suggested synthesis sentence to convert into prose:
\begin{quote}
The interpretability analyses therefore support a conservative reading of the classifier: BERT learned broad valence and common emotion categories, but rare-label performance and blended-affect errors limit its suitability for exact emotion-set prediction.
\end{quote}

\subsection{External Validation on Stories}
\label{subs:external-validation-stories}

\begin{itemize}
    \item Open by saying that Stories is the key out-of-domain, human-labeled test of the Lemotif-trained emotion model.
    \item Report the Stories test size: 114 reflections.
    \item Report trained BERT emotion performance on Stories: micro-$F_1 = .347$, macro-$F_1 = .246$, exact subset accuracy $= .018$, and hamming loss $= .167$.
    \item Report trained BERT derived sentiment performance on Stories: accuracy $= .544$ and macro-$F_1 = .490$.
    \item Compare trained and untrained BERT: emotion micro-$F_1$ .347 vs. .254; emotion macro-$F_1$ .246 vs. .038; sentiment accuracy .544 vs. .140.
    \item Interpret this comparison as evidence that Lemotif supervision matters, even though generalization remains limited.
    \item Report Gemma as a separate sentiment-only Stories comparison: sentiment accuracy $= .798$, macro-$F_1 = .674$, and average latency $= 20.50$ seconds per example.
    \item Clarify that the Gemma Stories comparison evaluates coarse sentiment, not the full 18-label Lemotif emotion taxonomy. Gemma models are open-weight language models based on Gemini-family research \parencite{gemmateam2024gemma}.
    \item Interpret the overall pattern as domain shift: Stories differs from Lemotif in length, context, tone, source population, and sentiment distribution.
\end{itemize}

\begin{table}[htbp]
    \centering
    \caption{External Stories test performance.}
    \label{tab:stories-external}
    \begin{tabular}{lrrrr}
        \toprule
        Model/task & Emotion micro-$F_1$ & Emotion macro-$F_1$ & Sent. acc. & Sent. macro-$F_1$ \\
        \midrule
        Trained BERT emotion model & .347 & .246 & .544 & .490 \\
        Untrained BERT control & .254 & .038 & .140 & .097 \\
        Gemma sentiment only & -- & -- & .798 & .674 \\
        \bottomrule
    \end{tabular}
\end{table}

\begin{figure}[htbp]
    \centering
    \includegraphics[width=\textwidth]{figures/results/main_text/fig07_stories_external_model_drop.png}
    \caption{Performance comparison between the in-domain Lemotif held-out test and the external Stories test. The drop in performance indicates substantial domain shift.}
    \label{fig:stories-external-model-drop}
\end{figure}

\begin{figure}[htbp]
    \centering
    \includegraphics[width=\textwidth]{figures/results/main_text/fig08_stories_per_emotion_f1.png}
    \caption{Per-emotion $F_1$ scores for the external Stories evaluation. The model generalizes better for frequent or affectively clear categories than for sparse or ambiguous labels.}
    \label{fig:stories-per-emotion-f1}
\end{figure}

\begin{figure}[htbp]
    \centering
    \includegraphics[width=\textwidth]{figures/results/main_text/fig09_stories_sentiment_model_comparison.png}
    \caption{Sentiment comparison on Stories for trained BERT, untrained BERT, and Gemma. Gemma is evaluated as a sentiment-only model and should not be interpreted as an 18-label emotion classifier in this comparison.}
    \label{fig:stories-sentiment-model-comparison}
\end{figure}

\subsection{WordPress Public-Journal Descriptive Analysis}
\label{subs:wordpress-public-journal-descriptive-analysis}

\begin{itemize}
    \item Start with the caveat: WordPress posts do not have human Lemotif labels, so this section cannot evaluate accuracy, $F_1$, or true model bias.
    \item State that the section analyzes model outputs descriptively to check whether inferred labels are plausible and whether the model simply reproduces the positive Lemotif training distribution.
    \item Report the WordPress scope: 72 requested sites, 72 successful sites, 10,671 analyzed reflections, 365-day window, and 365-reflection cap per site.
    \item Report the top inferred theme: God (11.0\%).
    \item Report the top inferred emotion: Anxious (48.9\%).
    \item State that public-journal predictions are not overwhelmingly positive; the strongest inferred emotions include Anxious, Frustrated, Surprised, Sad, Angry, Confused, and Happy.
    \item Use the theme-emotion heatmap to explain that Health, Sleep, and Work are especially associated with anxious/frustrated outputs, while God and Food show relatively more positive activation.
    \item Interpret possible skews cautiously: crawler/tag selection, public self-presentation, domain shift, model transfer bias, and higher affect density in negative entries.
    \item End with the RQ2 answer: WordPress inference does not formally validate the model, but it descriptively suggests that the model does not simply copy Lemotif's positive-class dominance.
\end{itemize}

\begin{figure}[htbp]
    \centering
    \includegraphics[width=\textwidth]{figures/results/main_text/fig10_wordpress_theme_emotion_heatmap.png}
    \caption{Inferred theme-emotion conditional prevalence in the public WordPress journal corpus. Because the corpus is not gold-labeled, this figure is interpreted descriptively rather than as model evaluation.}
    \label{fig:wordpress-theme-emotion-heatmap}
\end{figure}

% Appendix links to cite in prose:
% Appendix D: WordPress emotion prevalence, correlation matrix, and co-occurrence network.

\subsection{Results Summary by Research Question}
\label{subs:results-summary-rq}

\begin{itemize}
    \item RQ1, in-domain: BERT learned the Lemotif label space moderately well at the multi-label level and strongly at the derived sentiment level.
    \item RQ1, external Stories: performance dropped substantially, indicating that Lemotif-trained emotion classification does not yet generalize robustly to Stories reflections.
    \item RQ1, model comparison: Optuna improved micro-$F_1$, fixed inference improved label-density readability and sentiment accuracy, and Gemma improved Stories sentiment but was not a complete replacement for emotion taxonomy classification.
    \item RQ2, WordPress: the descriptive WordPress analysis suggests non-collapse into positive predictions, but no formal validation or bias metric is possible without gold labels.
    \item Bridge to Discussion: sentiment-level summarization is the stronger near-term use case, while exact multi-label emotion prediction requires more Stories-specific annotations and stronger rare-label handling.
\end{itemize}

Suggested closing paragraph to convert into final prose:
\begin{quote}
Overall, the results support a cautious interpretation of automatic affect analysis for reflective writing. The Lemotif-trained BERT model captured broad affective patterns and performed especially well when fine-grained emotions were collapsed into derived sentiment. However, the external Stories evaluation showed that exact Lemotif-style emotion classification remained sensitive to domain shift, label sparsity, and multi-label ambiguity. The WordPress analysis extended the pipeline to a larger public-journal setting, but only as descriptive evidence because the corpus was not human-labeled.
\end{quote}
